\documentclass[10pt,a4paper]{amsart}
\usepackage[margin=19mm]{geometry}
\usepackage{amsmath,amssymb,mathtools}
\usepackage[scr=rsfs,cal=euler]{mathalfa}
\usepackage{xfrac}
\usepackage[unicode,hidelinks,bookmarks=false]{hyperref}
\newcommand{\ie}{i.e.\ }
\newcommand{\cf}{cf.\ }
\newcommand{\resp}{respectively\ }
\newcommand{\Subsection}{\subsection}
\providecommand{\nobreakdash}{\nobreak\hbox{-}}
\setlength{\emergencystretch}{2em}
\multlinegap0pt
\usepackage{amssymb}
\usepackage{xcolor}
\usepackage{tikz}
\newtheorem{lemma}{Lemma}
\renewcommand{\le}{\leqslant}
\title{Decomposing tournaments into comparability graphs}
\author{Pierre Aboulker, Logan Crew, Julien Duron, Xinyue Fan, Hugo Jacob, R\'emy Kimbrough, Hidde Koerts, Benjamin Moore, Sophie Spirkl, St\'ephan Thomass\'e}
\date{Source: arXiv:2606.07748v1 (CC BY 4.0)}
\begin{document}
\maketitle

\noindent\emph{Source and licence.} Decomposing tournaments into comparability graphs. Version arXiv:2606.07748v1 (CC BY 4.0), distributed by arXiv under the Creative Commons Attribution 4.0 International licence, http://creativecommons.org/licenses/by/4.0/, as recorded in the arXiv licence metadata for this exact version and shown on the abstract page https://arxiv.org/abs/2606.07748v1. Full licence notice: \texttt{licenses/arxiv-2606.07748-CC-BY-4.0.txt}.\par\smallskip

\noindent\textit{Editorial scope.} Counted unit: the article's lemma lem:compgraphsubs (source0.tex lines 283--286). The article's complete original proof of it is delivered with it (source0.tex lines 288--300, one \texttt{proof} environment of 13 lines). No mathematics is written, repaired or re-derived: the statement and the proof are the article's own lines, in the article's own notation, and no retained line is substituted or re-flowed.  No further source-local context is needed: every cross-reference of every retained line resolves inside the two delivered spans.  Nothing was omitted from any retained span, so the package carries no omission record.  No retained line contains a citation, so the wrapper carries no bibliography and no bibliography entry.  No retained line contains a figure environment, an image-inclusion command or drawing code, so no image is delivered and the package declares no asset dependency.  Editorial formatting only, in addition: an AMS \texttt{amsart} preamble that restates the article's own declarations and macros used by the retained lines, namely the environments \texttt{lemma} on separate counters, together with the standard packages those lines need.  The retained environments print in this document as Lemma 1 in order of appearance, whereas the article prints the same items with the numbering of its own full document.  The complete native source of the article as distributed in the arXiv e-print for this exact version is bundled unchanged as \texttt{sources/202270/source0.tex}.
\begin{lemma}\label{lem:compgraphsubs}
Let $\mathcal{C}$ be a class of tournaments such that every $T\in\mathcal{C}$ satisfies $pod(T)\le m$.
Then every tournament $T\in\mathcal{C}^{subst}$ also satisfies $pod(T)\le m$.
\end{lemma}

\begin{proof}
    It suffices to show that if $T$ and $T'$ are both tournaments such that $pod(T) \leq m$ and $pod(T') \leq m$, then the same holds for all substitutions of $T'$ for a vertex of $T$. 

    Consider a substitution of $T'$ for a vertex $v$ of $T$ and let the resulting tournament be $T_{v \rightsquigarrow T'}$.  Let $P_1, \dots, P_m$ and $P'_1, \dots, P'_m$ be partial order decompositions of respectively $T$ and $T'$. We set $Q_i$ to be the spanning subdigraph of $T_{v \rightsquigarrow T'}$ obtained from substituting $P'_i$ for vertex $v$ in $P_i$. More formally: 
    \begin{align*}
    A(Q_i) ={}& A(P'_i) \cup \{e \in A(P_i) \mid v \notin e\}
 \cup \{(x,y) : x \in V(T'),\ (v,y) \in A(P_i)\}  \cup \{(y,x) : x \in V(T'),\ (y,v) \in A(P_i)\}.
\end{align*}

    Clearly $A(T_{v \rightsquigarrow T'}) = A(Q_1) \cup \dots \cup A(Q_m)$. It remains to show that each $Q_i$ is transitive. Let $a,b,c \in V(T_{v \rightsquigarrow T'})$ be three vertices such that $a \rightarrow_{T_{v \rightsquigarrow T'}} b$ and $b \rightarrow_{T_{v \rightsquigarrow T'}} c$. If $a, b, c \in V(T')$, then $(a,b), (b,c) \in A(P'_i)$, and $P'_i$ is transitive, so $(a,c) \in A(P'_i) \subseteq A(Q_i)$. If $|\{a, b, c\} \cap V(T')| \leq 1$, then the arcs corresponding to $(a, b)$ and $(b, c)$ in $P_i$ guarantee the existence of $(a,c)$ because $P_i$ is transitive. 

    Finally, observe that if $|\{a,b,c\} \cap V(T')| = 2$ then $b \in V(T')$ by definition of the substitution. Indeed, we have $x \Rightarrow V(T')$ for every $x \in V(T_{v \rightsquigarrow T'}) \setminus V(T')$. Then, exactly one of $a$ and $c$ is in $V(T')$ and in either case we have $a \to_{T_{v \rightsquigarrow T'}} c$ by definition of the substitution.
\end{proof}

\end{document}
